\section{Benchmarks and Evaluation Protocol}
\label{sec:data}

The paper's main scope is intentionally narrow. RWTD is the main natural benchmark and remains evaluation-only. STLD is the controlled synthetic companion with a canonical foreground side. Two additional routes remain in the appendix: a ControlNet-stitched bridge benchmark with exact binary supervision and a CAID breadth check. Neither route drives the main claim.

\paragraph{RWTD.} RWTD is the paper's anchor benchmark because it concentrates the failure modes that motivate recoverability: weak transitions, disconnected same-texture regions, and repeated-pattern distractors. Main results use the official evaluator plus a matched shared subset for direct comparison to the public TextureSAM rerun.

\paragraph{STLD.} STLD is the main synthetic companion. It is useful here because the shaped foreground side is canonical and exact supervision is available. That makes it a cleaner test of what the generated mask bank already reveals when evaluator ambiguity is minimal.

\paragraph{Supporting routes.} The appendix keeps two scope-limited routes. The ControlNet bridge benchmark is a controlled binary partition route generated by a ControlNet-conditioned stitched-regeneration pipeline; it preserves exact labels but uses more naturalistic transitions than classical hard mosaics. We use it as supporting evidence rather than a headline pillar because this paper does not establish a causal bridge-training effect. CAID is included as a domain-shift breadth check because its shoreline focus is less central to the paper's texture claim.

\paragraph{Metrics and subset conventions.} RWTD uses official mIoU and ARI under complementary binary partitions. STLD uses direct foreground mIoU and ARI because the shaped foreground side is canonical. The appendix bridge and CAID routes use partition-invariant mIoU and ARI. Every headline table keeps evaluator and subset conventions fixed within each comparison block; unmatched legacy numbers are not mixed into the main results.
